\documentclass[degree=master,degreetype=academic]{suepthesis}

\ctitle{上海电力大学论文格式测试标准文档}
\etitle{Standard Format Testing Document for SUEP Thesis}
\cauthor{张三丰}
\csupervisor{张无忌~教授}
\cmajor{软件工程}
\cdegree{学术硕士}
\cschool{计算机科学与技术学院}

% Required for dummy images and tables
\usepackage{graphicx}
\usepackage{booktabs}
\usepackage{caption}

% Figure and Table caption formatting based on the rules:
% Masters: Chinese only, Songti, Size 5 (wuhao), centered.
% We configure this globally.
\captionsetup[figure]{font=small, labelsep=space} % small is nominally 9pt / 10pt depending on class, let's explicit wuhao.
\DeclareCaptionFont{wuhaosong}{\wuhao\songti}
\captionsetup{font=wuhaosong, labelsep=quad}

\begin{document}

% 封面页
\MakeTitle

% 授权页
\makeauthpage

% ====================
% 中文摘要页
% ====================
\begin{cnabstract}
这是一篇用于严格测试“上海电力大学（SUEP）硕博士学位论文”排版规范的样例文档。本摘要部分的文字应为“宋体、小四号”，行距应该为“1.25倍”。而且每一个段落的开头应该留有两个汉字的缩进。

此外，摘要下方的“关键词”三字应该使用“黑体、四号”，而其后的具体关键词应该是“宋体、四号”，各关键词之间必须用逗号隔开，并且最后一个词汇之后不能有任何标点符号。

\cnkeywords{格式测试, 排版规范, LaTeX模板, 间距校验}
\end{cnabstract}

% ====================
% 英文摘要页
% ====================
\begin{enabstract}
This is a sample document used to strictly test the typesetting norms of the "Shanghai University of Electric Power (SUEP) Master's and Doctoral Dissertations". The text in this abstract section should be "Times New Roman, Size Small 4", and the line spacing should be "1.25 times". 

In addition, the word "KEY WORDS" below the abstract should use "Times New Roman, Size 4, Bold", followed by lower-case keywords separated by commas without a trailing punctuation mark.

\enkeywords{format testing, typographical norms, latex template, spacing verification}
\end{enabstract}

% ====================
% 目录页
% ====================
\tableofcontents

\clearpage
% ====================
% 正文开始
% ====================
\pagenumbering{arabic}
\pagestyle{suep_header}

\section{各级标题与段落格式测试}
本章主要用于测试一级、二级、三级标题的格式规范，以及普通正文的段落排版。

根据规范要求，正文部分的各种中文基本都应该是宋体、小四号，而数字和英文应该是 Times New Roman、小四号体。整篇文章采用首行缩进两个汉字，采用 1.25 倍的行距，并且段前段后距均为 0 磅。

所有的单数页页眉为当前章节标题，而偶数页页眉固定为“上海电力大学硕士/博士学位论文”。我们可以通过翻到下一页来验证偶数页页眉。

\subsection{二级标题的规范检查}
这是二级标题。按照要求，二级标题应该左对齐顶格，小三号宋体并且加黑，1.25倍行距，段前空一行，段后空 0.5 行。

我们在此处多写几段文字，以便将内容推送到第二页验证偶数页页眉是否正常显示。

这是另起的一段。所有的行文必须满足左右两端对齐的要求。这一页是奇数页，如果排版准确，页眉处的文字应该是“第一章 各级标题与段落格式测试”。

\subsubsection{三级标题测试}
这部分是三级标题的测试区。按照要求，三级标题应当左起空两个汉字字符，使用四号宋体并且加黑。段前应空 0.5 行，段后空 0 行。请仔细检查本标题左侧是否有相当于两个汉字宽度的留白。

\section{公式、图表排版测试}
本章我们将插入一些常见的数学公式、图和表，用以测试这些浮动体环境以及公式对齐是否符合标准。

\subsection{公式环境}
规范中要求“公式居中，序号顶右”。我们可以通过下方的方程测试此项规则：

\begin{equation}
  E = m c^2
\end{equation}

另一个多行的推导公式：
\begin{eqnarray}
  a &=& b + c \\
  c &=& d + e
\end{eqnarray}

\subsection{公式说明环境测试}
使用 \texttt{eqexpl} 环境来测试“式中：”的排版：

\begin{equation}
    PV = nRT
\end{equation}

\begin{eqexpl}
    \item[P] 压力 (Pressure)
    \item[V] 体积 (Volume)
    \item[n] 物质的量 (Amount of substance)
    \item[R] 理想气体常数 (Ideal gas constant)
    \item[T] 绝对温度 (Absolute temperature)
\end{eqexpl}

\subsection{图表环境}
对于硕士生来说，插图的图题只需使用中文，字号为“五号、宋体、居中排写”。我们将在这里插入一个虚拟图（图 \ref{fig:test}）用以检视图题的外观与文字。

\begin{figure}[H]
  \centering
  \rule{6cm}{3cm} % A simple black box acting as an image
  \bicaption{这是一个测试用途的虚拟插图}{This is a dummy illustration for testing purposes}
  \label{fig:test}
\end{figure}

表题也要遵循同样的规则，不过表格一般推荐使用“三线表”格式。如下表 \ref{tab:test} 所示：

\begin{table}[H]
  \centering
  \caption{这是一个测试使用的数据表格}
  \label{tab:test}
  \begin{tabular}{ccc}
    \toprule
    测试项目 & 期望字号 & 期望字体 \\
    \midrule
    图题/表题 & 五号 & 宋体 \\
    正文内容 & 小四号 & 宋体/Times New Roman \\
    \bottomrule
  \end{tabular}
\end{table}

\input{chapters/chap03}

\section{参考文献}
为了测试参考文献目录，我们在文末通常会有相关的 `thebibliography` 环境配置。由于尚未提供完整的参考文献环境规范，我们这里的验证重心停留在正文及标题格式。

\end{document}
